\documentclass[11pt]{article}

\usepackage{amsmath,amssymb,amsthm,mathtools}
\usepackage[margin=1in]{geometry}
\usepackage{microtype}

\newtheorem{theorem}{Theorem}[section]
\newtheorem{lemma}[theorem]{Lemma}
\newtheorem{corollary}[theorem]{Corollary}
\newtheorem{remark}[theorem]{Remark}

\newcommand{\dd}{\,\mathrm d}

\title{Triangle Books with a Page-Rooted Triangle--Leaf Branch}
\author{}
\date{}

\begin{document}

\maketitle

\begin{abstract}
Let $B_m^{\rm new}$ and $B_m^{\rm page}$ be obtained from a book of
$m\geq1$ triangles on a common spine by attaching a new triangle along the
edge from one fixed spine endpoint to one exceptional page.  Attach one
leaf either to the new triangle vertex or to the exceptional page,
respectively.  We prove that every graphon of edge density
$p\in[1/2,1]$ satisfies, for either graph,
\[
 t(B_m^\star,W)\geq p^2(2p-1)^{m+1}
 =\Phi_{B_m^\star}(p).
\]
The proof symmetrizes the exceptional branch over the book pages.  Two
fractional edge-rooted estimates follow from edge Jensen and the
real-exponent weighted rooted-triangle inequality; one of them also uses a
H\"older compression.  The cases $m=2$ prove NetworkX Atlas graphs $137$
and $139$.
\end{abstract}


%%%%%%%%%%%%%%%%%%%%%%%%%%%%%%%%%%%%%%%%%%%%%%%%%%%%%%%%%%%%%%%%%%%%%%%%%%%%%%%
\section{The graph family and its target}
%%%%%%%%%%%%%%%%%%%%%%%%%%%%%%%%%%%%%%%%%%%%%%%%%%%%%%%%%%%%%%%%%%%%%%%%%%%%%%%

Let $(\Omega,\mathcal A,\mu)$ be an arbitrary probability space and let
$W\colon\Omega^2\to[0,1]$ be symmetric and jointly measurable.  Put
\[
 p=\int_{\Omega^2}W(x,y)\dd\mu(x)\dd\mu(y),
 \qquad
 d(x)=\int_\Omega W(x,y)\dd\mu(y).
\]

For $m\geq1$, define two graphs as follows.  Their spine vertices are
$a,b$, their page vertices are $z_1,\ldots,z_m$, and their two additional
vertices are $u,v$.  Their common edges are
\[
 ab,\qquad az_i,bz_i\quad(1\leq i\leq m),
 \qquad au,z_1u,
\]
together with $uv$ for $B_m^{\rm new}$ and with $z_1v$ for
$B_m^{\rm page}$.  Thus $a,z_1,u$ is the added triangle; the superscript
records whether $v$ is a leaf at its new vertex $u$ or at the old page
vertex $z_1$.

After the ordered spine has been properly coloured, each book page and the
new triangle vertex have $x-2$ choices, while the leaf has $x-1$ choices.
Consequently
\begin{equation}
 \chi_{B_m^\star}(x)=x(x-1)^2(x-2)^{m+1},
 \qquad \chi(B_m^\star)=3.
 \label{eq:chromatic}
\end{equation}
For $q=1-p>0$, substitution in the chromatic target gives
\begin{equation}
 \Phi_{B_m^\star}(p)
 =q^{m+4}\chi_{B_m^\star}(1/q)
 =p^2(2p-1)^{m+1}.
 \label{eq:target}
\end{equation}
The right side is polynomial and supplies the continued value at $p=1$.


%%%%%%%%%%%%%%%%%%%%%%%%%%%%%%%%%%%%%%%%%%%%%%%%%%%%%%%%%%%%%%%%%%%%%%%%%%%%%%%
\section{Real weighted triangles}
%%%%%%%%%%%%%%%%%%%%%%%%%%%%%%%%%%%%%%%%%%%%%%%%%%%%%%%%%%%%%%%%%%%%%%%%%%%%%%%

Define the rooted triangle density
\[
 \tau(x)=\int_{\Omega^2}
 W(x,y)W(x,z)W(y,z)\dd\mu(y)\dd\mu(z)
\]
and, for every real $s\geq0$,
\[
 M_s=\int_\Omega d(x)^s\dd\mu(x),
\]
with $d^0=1$ also on the zero-degree set.

\begin{lemma}[Weighted rooted triangle]
\label{lem:weighted-triangle}
If $p>0$, then for every real $h\geq0$,
\begin{equation}
 \int_\Omega d(x)^h\tau(x)\dd\mu(x)
 \geq\frac{2p-1}{p}M_{h+2}.
 \label{eq:weighted-triangle}
\end{equation}
\end{lemma}

\begin{proof}
For a bounded kernel $K$, write
$(T_Kf)(x)=\int K(x,y)f(y)\dd\mu(y)$ and use the usual $L^2$ inner
product.  Set
\[
 I_h=\int d^h\tau,
 \qquad J_h=\langle d^h,T_Wd\rangle.
\]
The inequality $rs\geq r+s-1$ on $[0,1]^2$, applied to the two edges
from the root of the triangle and multiplied by $d(x)^hW(y,z)$, yields
\begin{equation}
 I_h\geq2J_h-pM_h.
 \label{eq:first-reduction}
\end{equation}

Put $U=1-W$, $u=T_U\mathbf1=1-d$, and $q=1-p$.  Since
\[
 T_Wd=d-q\mathbf1+T_Uu,
\]
we have
\begin{equation}
 J_h=M_{h+1}-qM_h+\langle d^h,T_Uu\rangle.
 \label{eq:J}
\end{equation}
Symmetry of $U$ gives
\begin{align*}
 &2\left(\langle d^h,T_Uu\rangle-\int d^hu^2\right)\\
 &\quad=
 \int_{\Omega^2}U(x,y)
 \bigl(u(y)-u(x)\bigr)
 \bigl(d(x)^h-d(y)^h\bigr)\dd\mu^2.
\end{align*}
The integrand is nonnegative because $d=1-u$ and $t\mapsto t^h$ is
nondecreasing.  Therefore
\[
 \langle d^h,T_Uu\rangle
 \geq M_h-2M_{h+1}+M_{h+2}.
\]
Substitution into \eqref{eq:J} and \eqref{eq:first-reduction} gives
\[
 I_h\geq2M_{h+2}-2M_{h+1}+pM_h.
\]
After subtracting the right side of \eqref{eq:weighted-triangle}, the
remaining quantity is exactly
\[
 \frac1p\int_\Omega d(x)^h(d(x)-p)^2\dd\mu(x)\geq0.
\]
No integrality of $h$ was used.
\end{proof}


%%%%%%%%%%%%%%%%%%%%%%%%%%%%%%%%%%%%%%%%%%%%%%%%%%%%%%%%%%%%%%%%%%%%%%%%%%%%%%%
\section{A fractional edge-rooted paw estimate}
%%%%%%%%%%%%%%%%%%%%%%%%%%%%%%%%%%%%%%%%%%%%%%%%%%%%%%%%%%%%%%%%%%%%%%%%%%%%%%%

For $x,y\in\Omega$ and real $s\geq0$, put
\[
 H_s(x,y)=\int_\Omega W(x,z)W(y,z)d(z)^s\dd\mu(z).
\]
Thus $H_0$ is the common-neighbour kernel, while $H_1$ is the rooted
density obtained by putting one leaf at the new common neighbour.

\begin{lemma}[Fractional edge-rooted paw]
\label{lem:fractional-paw}
Let $p\in[1/2,1]$ and $0<\alpha\leq1$.  Then
\begin{equation}
 \int_{\Omega^2}W(x,y)H_0(x,y)H_1(x,y)^\alpha\dd\mu^2
 \geq \bigl(p(2p-1)\bigr)^{1+\alpha}.
 \label{eq:fractional-paw}
\end{equation}
\end{lemma}

\begin{proof}
Put $s=\alpha/(1+\alpha)$.  H\"older's inequality applied to the measure
$W(x,z)W(y,z)\dd\mu(z)$ gives, pointwise in $(x,y)$,
\[
 H_s\leq H_0^{1-s}H_1^s.
\]
Since $(1-s)(1+\alpha)=1$ and $s(1+\alpha)=\alpha$, this is equivalent to
\begin{equation}
 H_0H_1^\alpha\geq H_s^{1+\alpha}.
 \label{eq:holder-compression}
\end{equation}

Because $p\geq1/2>0$, the measure $W(x,y)\dd\mu^2/p$ is a probability
measure.  Jensen's inequality and \eqref{eq:holder-compression} imply
\begin{align}
 \int W H_0H_1^\alpha
 &\geq\int W H_s^{1+\alpha}\notag\\
 &\geq p^{-\alpha}\left(\int W H_s\right)^{1+\alpha}.
 \label{eq:edge-jensen}
\end{align}
Fubini gives the exact identity
\[
 \int_{\Omega^2}W(x,y)H_s(x,y)\dd\mu^2
 =\int_\Omega d(z)^s\tau(z)\dd\mu(z).
\]
By Lemma~\ref{lem:weighted-triangle} and Jensen for the convex power
$t\mapsto t^{s+2}$,
\[
 \int W H_s
 \geq\frac{2p-1}{p}M_{s+2}
 \geq p^{s+1}(2p-1).
\]
Insert this in \eqref{eq:edge-jensen}.  The choice of $s$ gives
\[
 -\alpha+(s+1)(1+\alpha)=1+\alpha,
\]
and hence the claimed lower bound.
\end{proof}


%%%%%%%%%%%%%%%%%%%%%%%%%%%%%%%%%%%%%%%%%%%%%%%%%%%%%%%%%%%%%%%%%%%%%%%%%%%%%%%
\section{A fractional exceptional-page estimate}
%%%%%%%%%%%%%%%%%%%%%%%%%%%%%%%%%%%%%%%%%%%%%%%%%%%%%%%%%%%%%%%%%%%%%%%%%%%%%%%

\begin{lemma}[Fractional exceptional page]
\label{lem:fractional-page}
Let $p\in[1/2,1]$ and $0<\alpha\leq1$.  Then
\begin{equation}
 \int_{\Omega^2}W(x,y)d(y)^\alpha H_0(x,y)^{1+\alpha}\dd\mu^2
 \geq \bigl(p(2p-1)\bigr)^{1+\alpha}.
 \label{eq:fractional-page}
\end{equation}
\end{lemma}

\begin{proof}
Put $s=\alpha/(1+\alpha)$.  Since
\[
 d(y)^\alpha H_0(x,y)^{1+\alpha}
 =\bigl(d(y)^sH_0(x,y)\bigr)^{1+\alpha},
\]
Jensen under the edge-biased probability measure gives
\[
 \int Wd(y)^\alpha H_0^{1+\alpha}
 \geq p^{-\alpha}
 \left(\int W(x,y)d(y)^sH_0(x,y)\dd\mu^2\right)^{1+\alpha}.
\]
Fubini gives
\[
 \int_{\Omega^2}W(x,y)d(y)^sH_0(x,y)\dd\mu^2
 =\int_\Omega d(y)^s\tau(y)\dd\mu(y).
\]
The final two lines of the proof of Lemma~\ref{lem:fractional-paw} now
apply verbatim and give \eqref{eq:fractional-page}.
\end{proof}


%%%%%%%%%%%%%%%%%%%%%%%%%%%%%%%%%%%%%%%%%%%%%%%%%%%%%%%%%%%%%%%%%%%%%%%%%%%%%%%
\section{The two branch orientations}
%%%%%%%%%%%%%%%%%%%%%%%%%%%%%%%%%%%%%%%%%%%%%%%%%%%%%%%%%%%%%%%%%%%%%%%%%%%%%%%

\begin{theorem}[Triangle books with one page-rooted triangle--leaf branch]
\label{thm:main}
For every integer $m\geq1$, every graphon $W$ of edge density
$p\in[1/2,1]$, and $\star\in\{\mathrm{new},\mathrm{page}\}$,
\[
 \boxed{
 t(B_m^\star,W)\geq p^2(2p-1)^{m+1}
 =\Phi_{B_m^\star}(p).
 }
\]
This is the complete interval required by $\chi(B_m^\star)=3$.
\end{theorem}

\begin{proof}
First take $B_m^{\rm new}$.  Integrating the added leaf and then its
neighbour gives the factor $H_1(x,z_1)$ on the exceptional book page.  The
$m$ page variables form one orbit in the underlying book.  Average the $m$
equal integral representations obtained by moving the exceptional factor
among the pages, and apply arithmetic--geometric mean pointwise.  With
$x,y$ denoting the spine images, this gives
\begin{equation}
 t(B_m^{\rm new},W)
 \geq\int_{\Omega^2}W(x,y)R(x,y)^m\dd\mu^2,
 \label{eq:page-factorization}
\end{equation}
where
\[
 R(x,y)=\int_\Omega
 W(x,z)W(y,z)H_1(x,z)^{1/m}\dd\mu(z).
\]
The pages factor independently after symmetrization, which justifies
\eqref{eq:page-factorization}.

Jensen under the edge-biased probability measure gives
\begin{equation}
 t(B_m^{\rm new},W)
 \geq p^{1-m}\left(\int_{\Omega^2}W(x,y)R(x,y)\dd\mu^2\right)^m.
 \label{eq:outer-jensen}
\end{equation}
Fubini and the definition of $H_0$ show that the inner integral is
\[
 \int_{\Omega^2}W(x,z)H_0(x,z)H_1(x,z)^{1/m}\dd\mu(x)\dd\mu(z).
\]
Lemma~\ref{lem:fractional-paw} with $\alpha=1/m$ bounds this by
\[
 \bigl(p(2p-1)\bigr)^{1+1/m}.
\]
Substitution into \eqref{eq:outer-jensen} yields
\[
 t(B_m^{\rm new},W)
 \geq p^{1-m}\bigl(p(2p-1)\bigr)^{m+1}
 =p^2(2p-1)^{m+1},
\]
which is \eqref{eq:target}.

For $B_m^{\rm page}$, integrating the leaf at $z_1$ and the new triangle
vertex $u$ gives the exceptional factor $d(z_1)H_0(x,z_1)$.  The identical
page-orbit symmetrization and outer edge Jensen give
\[
 t(B_m^{\rm page},W)
 \geq p^{1-m}\left(\int_{\Omega^2}W(x,y)S(x,y)\dd\mu^2\right)^m,
\]
where
\[
 S(x,y)=\int_\Omega W(x,z)W(y,z)
 \bigl(d(z)H_0(x,z)\bigr)^{1/m}\dd\mu(z).
\]
Fubini turns the integral in parentheses into
\[
 \int_{\Omega^2}W(x,z)d(z)^{1/m}H_0(x,z)^{1+1/m}
 \dd\mu(x)\dd\mu(z).
\]
Lemma~\ref{lem:fractional-page} with $\alpha=1/m$ bounds this by
$\bigl(p(2p-1)\bigr)^{1+1/m}$.  The same exponent calculation therefore
proves the asserted bound for $B_m^{\rm page}$.
\end{proof}

At $p=1/2$ the target is zero and all divisions above are only by the
positive number $p$.  At $p=1$, $W=1$ almost everywhere and both sides are
one.  For every balanced complete $k$-partite graphon, $k\geq2$, one has
$d=p$ and $H_s=p^s(2p-1)$ on every supported edge; hence all H\"older,
Jensen, and symmetrization steps are equalities, as required.


%%%%%%%%%%%%%%%%%%%%%%%%%%%%%%%%%%%%%%%%%%%%%%%%%%%%%%%%%%%%%%%%%%%%%%%%%%%%%%%
\section{Atlas 137, Atlas 139, and scope}
%%%%%%%%%%%%%%%%%%%%%%%%%%%%%%%%%%%%%%%%%%%%%%%%%%%%%%%%%%%%%%%%%%%%%%%%%%%%%%%

\begin{corollary}[Atlas 137 and 139]
NetworkX Atlas graphs $137$ and $139$, with graph6 codes \texttt{EjsG} and
\texttt{Eju?}, respectively, are positive.  Each has order six, size eight,
chromatic number three, and
\[
 \chi_{F_{139}}(x)=x(x-1)^2(x-2)^3.
\]
For $i\in\{137,139\}$, every graphon of edge density $p\in[1/2,1]$
satisfies
\[
 t(F_{139},W)\geq p^2(2p-1)^3=\Phi_{F_{139}}(p).
\]
\end{corollary}

\begin{proof}
For Atlas $137$, use spine $13$, ordinary page $2$, exceptional page $4$,
added triangle vertex $0$, and leaf $5$ at page $4$.  Its edge list is
\[
 01,04,12,13,14,23,34,45,
\]
so it is $B_2^{\rm page}$.  For Atlas $139$, use the same roles but put the
leaf at the new vertex $0$.  Its edge list is
\[
 01,04,05,12,13,14,23,34,
\]
so it is $B_2^{\rm new}$.  Theorem~\ref{thm:main} applies to both.
\end{proof}

All integrands in the proof are bounded, nonnegative, and measurable, so
Tonelli--Fubini applies directly on an arbitrary probability space.  The
argument concerns ordinary homomorphism densities; it uses no finite-step
reduction, injective-copy interpretation, regularity, atomlessness, or
minimum-degree hypothesis.  Zero degrees cause no quotient.

The theorem requires one triangle--leaf branch rooted at a page and one
fixed spine endpoint of a single triangle book, in either of the two stated
leaf orientations.  It does not assert arbitrary leaf,
triangle-page, edge-gluing, or mixed-chordal closure.  In particular, its
proof uses both page-orbit symmetrization and the specific edge-rooted
H\"older compression \eqref{eq:holder-compression}.


%%%%%%%%%%%%%%%%%%%%%%%%%%%%%%%%%%%%%%%%%%%%%%%%%%%%%%%%%%%%%%%%%%%%%%%%%%%%%%%
\appendix
\section{What the Lean formalization actually proves}
%%%%%%%%%%%%%%%%%%%%%%%%%%%%%%%%%%%%%%%%%%%%%%%%%%%%%%%%%%%%%%%%%%%%%%%%%%%%%%%

The Lean development formalizes the case $m=2$, which is the only case the
catalogue needs.  The modules are
\texttt{lean/Taeyoung/Methods/PagePawBranch/Core.lean} and
\texttt{Rows.lean}; the two catalogue theorems are
\texttt{satisfiesLowerBound\_bookNew} (Atlas $139$) and
\texttt{satisfiesLowerBound\_bookPage} (Atlas $137$).  The formalization
departs from the argument above at two places, and this appendix says where.


\subsection{The page-orbit symmetrization is a H\"older inequality}
\label{sec:app-symmetrization}

The proof of Theorem~\ref{thm:main} obtains \eqref{eq:page-factorization} by
averaging $m$ integral representations obtained by moving the exceptional
factor among the pages, and applying arithmetic--geometric mean pointwise.
Those $m$ representations are indeed equal, but for a trivial reason: the $m$
page factors are $m$ factors of one product, so naming which of them carries
the exceptional weight changes nothing, and the average is the same single
quantity.  Symmetrization therefore contributes no inequality, and the step
from that product to $\int\!\!\int WR^m$ is carried entirely by the pointwise
bound $R^m\leq GH_0^{m-1}$ below, which is H\"older, not
arithmetic--geometric mean.

Peeling the leaf and the glued vertex gives the exact identity
\begin{equation}
 t(B_m^{\rm new},W)
 =\int_{\Omega^2}W(x,y)\,H_0(x,y)^{m-1}\,G(x,y)\dd\mu^2,
 \qquad
 G(x,y)=\int_\Omega W(x,z)W(y,z)H_1(x,z)\dd\mu(z),
 \label{eq:app-peel}
\end{equation}
with no inequality yet used.  H\"older on the finite measure
$W(x,z)W(y,z)\dd\mu(z)$, applied to the pair of exponents $m$ and $m/(m-1)$,
gives
\begin{equation}
 R(x,y)=\int_\Omega W(x,z)W(y,z)H_1(x,z)^{1/m}\dd\mu(z)
 \leq G(x,y)^{1/m}H_0(x,y)^{(m-1)/m},
 \label{eq:app-holder}
\end{equation}
that is $R^m\leq GH_0^{m-1}$, which together with \eqref{eq:app-peel} is
exactly \eqref{eq:page-factorization}.  At $m=2$ the two exponents are both
$2$ and \eqref{eq:app-holder} is Cauchy--Schwarz.  The same remark applies
verbatim to $B_m^{\rm page}$, with $H_1(x,z)$ replaced by $d(z)H_0(x,z)$.

The Lean proof therefore never symmetrizes.  It works with the
\emph{branch operator}
\[
 \Lambda_h(x,y)=\int_\Omega W(x,z)W(y,z)h(x,z)\dd\mu(z),
\]
of which $H_0=\Lambda_1$ is the special case $h\equiv1$, proves
\eqref{eq:app-peel} in the form
\[
 t=\int_{\Omega^2}W\cdot H_0\cdot\Lambda_h,
 \qquad
 h=H_1\ \ (B_2^{\rm new}),
 \qquad
 h(x,z)=d(z)H_0(x,z)\ \ (B_2^{\rm page}),
\]
and replaces \eqref{eq:app-holder} by
\[
 \Lambda_{\sqrt h}(x,y)^2\leq H_0(x,y)\,\Lambda_h(x,y)
 \qquad\text{(\texttt{sq\_branchOp\_le}),}
\]
one application of the weighted Cauchy--Schwarz inequality
$(\int A\eta)^2\leq(\int A)(\int A\eta^2)$ with $A=W(x,\cdot)W(y,\cdot)$ and
$\eta=\sqrt{h(x,\cdot)}$.  This is the project's
\texttt{PureChordal.WeightedCauchySchwarz}, already used for the page
compression $H_{s/2}^2\leq H_0H_s$.


\subsection{Edge Jensen at exponent $3/2$, without convexity}
\label{sec:app-jensen}

Both fractional estimates of Sections~3 and~4 use Jensen for the convex power
$t\mapsto t^{1+\alpha}$ under the edge-biased probability measure
$W\dd\mu^2/p$.  At $\alpha=1/2$ the exponent is $3/2$, and the development has
no lemma for Jensen at a fractional power on a biased measure.  It is not
needed: writing
\[
 A=\int_{\Omega^2}WF,
 \qquad
 S=\int_{\Omega^2}W\sqrt F,
 \qquad
 T=\int_{\Omega^2}WF\sqrt F
\]
for a measurable $F$ with values in $[0,1]$, two applications of the same
weighted Cauchy--Schwarz inequality --- at weight $W$ with $\eta=\sqrt F$, and
at weight $W\sqrt F$ with the same $\eta$ --- give
\[
 S^2\leq pA,
 \qquad
 A^2\leq ST.
\]
Eliminating $S$ yields $A^4\leq S^2T^2\leq pAT^2$, hence
\begin{equation}
 A^3\leq pT^2
 \qquad\text{(\texttt{cube\_edge\_le}),}
 \label{eq:app-cube}
\end{equation}
which is the $\alpha=1/2$ case of \eqref{eq:edge-jensen} after clearing
denominators.  The case $A=0$ is separate and immediate.  So the whole
formalization uses one inequality, weighted Cauchy--Schwarz, at four different
weights, and no genuine H\"older inequality and no convexity lemma anywhere.


\subsection{The two comparison kernels}
\label{sec:app-kernels}

With \eqref{eq:app-cube} in hand, each row supplies a kernel $F$ with two
properties: the pointwise compression $F\sqrt F\leq\sqrt h\,H_0$, and the
scalar bound $\int\!\!\int WF\geq p^{4/3}(2p-1)$.  The generic conclusion
\[
 t\geq p^{-2}\Bigl(\int_{\Omega^2}WF\Bigr)^{3}\geq p^2(2p-1)^3
 \qquad\text{(\texttt{branch\_bound})}
\]
then holds for both rows at once.  The two choices are:
\begin{itemize}
 \item $B_2^{\rm new}$: $F=H_{1/3}$.  The compression squares to
 $H_{1/3}^3\leq H_0^2H_1$, which is \eqref{eq:holder-compression} at
 $\alpha=1/2$ and is already in the development as
 \texttt{Link.cube\_pageOp\_le}, itself two staggered Cauchy--Schwarz steps.
 The scalar bound is $\int\!\!\int WH_{1/3}=\int d^{1/3}\tau$
 (\texttt{Link.integral\_edge\_pageOp}) together with
 Lemma~\ref{lem:weighted-triangle} at $h=1/3$ and Jensen at exponent $7/3$.
 \item $B_2^{\rm page}$: $F(x,z)=d(z)^{1/3}H_0(x,z)$.  The compression is an
 \emph{identity}, both sides being $d(z)^{1/2}H_0(x,z)^{3/2}$.  The scalar
 bound needs the Fubini identity
 \[
  \int_{\Omega^2}W(x,y)\,d(y)^{s}H_0(x,y)\dd\mu^2=\int_\Omega d(z)^s\tau(z)\dd\mu(z),
 \]
 the counterpart of \texttt{integral\_edge\_pageOp} with the leaf weight
 carried by a spine endpoint rather than by the page vertex.  It is proved as
 \texttt{integral\_edge\_degree\_pageOp}.
\end{itemize}
In both cases the only fractional exponent surviving to the end is the
$p^{4/3}$ of the scalar bound, and it is cubed back to $p^4$ exactly as in the
three-page row \texttt{PageBook.Atlas138}.


\subsection{Scope of the formalization}
\label{sec:app-scope}

Only $m=2$ is formalized.  The general-$m$ theorem is not: the two ingredients
that would have to be generalized are \eqref{eq:app-holder} at $m$ factors and
\eqref{eq:app-cube} at exponent $1+1/m$, and both are stated in the Lean
development at the square root only.

The Lean graphs use their own labelling, chosen so that peeling produces
\eqref{eq:app-peel} directly: spine $0,1$; exceptional page $2$; ordinary page
$3$; glued vertex $4$ on the page edge $0$--$2$; leaf $5$, at the glued vertex
for $B_2^{\rm new}$ and at the exceptional page for $B_2^{\rm page}$.  The
generated Atlas modules \texttt{Examples/Graph139.lean} and
\texttt{Examples/Graph137.lean} carry the bound to the Atlas edge lists along
an explicit relabelling whose adjacency correspondence is checked by kernel
computation.  The chromatic data \eqref{eq:chromatic} is proved, not assumed,
by a three-step clique-attachment tower over $K_3$: pages attach to the
$2$-cliques $\{a,b\}$ and $\{a,z_1\}$, and the leaf to a single vertex, giving
$x(x-1)^2(x-2)^3$ with chromatic number three.

Both catalogue theorems, and the twenty-one supporting lemmas listed in
\texttt{Taeyoung/CheckVerified.lean}, depend on exactly
\texttt{[propext, Classical.choice, Quot.sound]}.

\end{document}
